\begin{description}
\item[(T9.1)]
Let apogee and perigee distances be $r_{\mathrm{a}}$ and $r_{\mathrm{p}}$
respectively.
\begin{align*}
r_{\mathrm{p}} &= R_\oplus + h^i_{\mathrm{p}} = (6371 + 264)\,\mathrm{km}
  = 6635\,\mathrm{km} \\
r_{\mathrm{a}} &= R_\oplus + h^i_{\mathrm{a}} = (6371 + 23904)\,\mathrm{km}
  = 30\,275\,\mathrm{km}
\end{align*}
Conservation of energy and angular momentum gives
\[ E = -\frac{GmM_\oplus}{r_{\mathrm{p}} + r_{\mathrm{a}}} \]
Total energy at perigee
\begin{align*}
\frac{1}{2}m{v_{\mathrm{p}}}^2 - \frac{GmM_\oplus}{r_{\mathrm{p}}}
  &= E = -\frac{GmM_\oplus}{r_{\mathrm{p}} + r_{\mathrm{a}}} \\
\frac{1}{2}{v_{\mathrm{p}}}^2
  &= \frac{GM_\oplus}{r_{\mathrm{p}}}
     \left(1 - \frac{r_{\mathrm{p}}}{r_{\mathrm{p}} + r_{\mathrm{a}}}\right) \\
\therefore v_{\mathrm{p}}
  &= \sqrt{2\,\frac{GM_\oplus}{r_{\mathrm{a}} + r_{\mathrm{p}}}\,
     \frac{r_{\mathrm{a}}}{r_{\mathrm{p}}}} \\
  &= \sqrt{\frac{2 \times 6.674 \times 10^{-11} \times 5.972 \times 10^{24}
     \times 3.0275 \times 10^{7}}
     {6.635 \times 10^{6} \times (3.0275 + 6.635 \times 10^{6})}} \\
% sic: the source prints the denominator as "(3.0275 + 6.635 x 10^6)";
% it should read (3.0275 x 10^7 + 6.635 x 10^6).
  &= 9.929\,\mathrm{km\,s^{-1}}
\end{align*}
As the engine burn is just 41.6\,s, we assume that the entire impulse is
applied instantaneously at perigee. The impulse is
$J = 1.73 \times 10^{5}\,\mathrm{kg\,m\,s^{-1}}$. Note that the total mass of
MOM must include the fuel, so we have to use
\[ m = 500 + 852 = 1352\,\mathrm{kg}. \]
Change in velocity due to impulse at perigee is
\[ \Delta v = \frac{J}{m} = \frac{1.73 \times 10^{5}}{1352}
   = 128.0\,\mathrm{m\,s^{-1}} \]
The new velocity will be given by (we use $'$ symbol to denote quantities
after the first orbit-raising maneuvre)
\[ v'_{\mathrm{p}} = v_{\mathrm{p}} + \Delta v = 10.06\,\mathrm{km\,s^{-1}} \]
The perigee remains unchanged. So we get $r'_{\mathrm{p}} = r_{\mathrm{p}}$.
Since the satellite is moving faster, the new apogee will be higher.
\begin{align*}
v'_{\mathrm{p}}
  &= \sqrt{2GM_\oplus \frac{r'_{\mathrm{a}}}
     {r'_{\mathrm{p}}(r'_{\mathrm{a}} + r'_{\mathrm{p}})}} \\
\therefore 1 + \frac{r'_{\mathrm{p}}}{r'_{\mathrm{a}}}
  &= \frac{2GM_\oplus}{(v'_{\mathrm{p}})^2 \times r'_{\mathrm{p}}} \\
  &= \frac{2 \times 6.674 \times 10^{-11} \times 5.972 \times 10^{24}}
     {(10.06 \times 10^{3})^2 \times 6.635 \times 10^{6}} = 1.188 \\
r'_{\mathrm{a}} &= \frac{6635}{0.188} = 35\,380\,\mathrm{km} \\
h_{\mathrm{a}} &= 35380 - 6371 \\
  &= \boxed{29\,009\,\mathrm{km}}
\end{align*}
Acceptable range: $\pm 150$\,km

\item[(T9.2)]
As seen above,
\begin{align*}
\frac{r'_{\mathrm{p}}}{r'_{\mathrm{a}}} &= 0.188 = \frac{1-e}{1+e} \\
\therefore e &= \frac{1 - 0.188}{1 + 0.188} = \boxed{0.683}
\end{align*}
Acceptable range: $\pm 0.002$

The new orbital semi-major axis and orbital period will be,
\begin{align*}
a' &= \frac{r'_{\mathrm{a}} + r'_{\mathrm{p}}}{2} \\
   &= \frac{35380 + 6635}{2} = 20\,933\,\mathrm{km} \\
P &= 2\pi\sqrt{\frac{a'^3}{GM_\oplus}} \\
  &= 2\pi\sqrt{\frac{(2.0933 \times 10^{7})^3}
     {6.674 \times 10^{-11} \times 5.972 \times 10^{24}}}
   = \boxed{30\,136\,\mathrm{s} = 8.37\,\mathrm{h}}
\end{align*}
Acceptable range: $\pm 0.1$\,h
\end{description}
